\section{Experimental Setup}

\subsection{Freeform device parameterization and physical stack}

The study focuses on freeform binary metasurfaces represented on a $64\times64$ planar grid. This representation is intentionally broad: it avoids restricting the search to a small analytical family and makes the inverse problem difficult in the way freeform operator design is actually difficult. The structures are binary in the sense relevant to fabrication-oriented patterning, and the data-generation routines impose symmetry priors consistent with the intended design space. In particular, the structure generation and augmentation steps enforce $C_4+\sigma_x+\sigma_y$ symmetry, or equivalently a $D_4$-type symmetry prior. This choice is not an arbitrary geometric restriction; it reduces obvious azimuthal redundancy, promotes in-plane isotropy, and makes the $0^{\circ}$--$45^{\circ}$ sector the representative fundamental wedge from which the full structure or angular response can be recovered by symmetry folding. The resulting design space is broad enough to host several operator families, but sparse enough that naive search rapidly becomes wasteful.

All electromagnetic labeling and verification in the present work are carried out with a fixed physical stack corresponding to a silicon patterned layer embedded between air input and silica output media. The nominal geometric and material parameters used throughout the code base are a $500\,\mathrm{nm}$ period, a $500\,\mathrm{nm}$ patterned-layer thickness, silicon as the structure material, air as the incident medium, and $\mathrm{SiO}_2$ as the output medium. These choices are held fixed not because they are universally optimal, but because the paper is concerned with the inverse-design formulation rather than with exhaustive material-stack exploration.

\subsection{Angle--frequency sampling and response representation}

The response field is sampled on a wavelength grid from $800$ to $1300\,\mathrm{nm}$ in $50\,\mathrm{nm}$ steps and an angular grid from $-40^{\circ}$ to $40^{\circ}$ in $5^{\circ}$ steps. The resulting $11\times17$ grid is dense enough to encode the broad transfer-function geometry targeted in this study while remaining computationally tractable for RCWA labeling and reranking. Two polarization channels are retained, corresponding to the $T_{pp}$ and $T_{ss}$ magnitude responses used in the training data and in the task definitions. Each sample is therefore represented by a pair of angle--wavelength maps rather than by a scalar score or a one-dimensional spectrum.

This discretization has practical consequences for what can and cannot be claimed. It supports operator families whose distinguishing physics is visible at the scale of the chosen angular and spectral window, including second-order and fourth-order differentiation, low-pass and high-pass angular filtering, polarization-independent operation, and polarization multiplexing. At the same time, it does not claim to resolve arbitrarily sharp resonant features or all possible phase-sensitive effects. The goal of the paper is to show that a unified response-space formulation is viable within a realistic computational budget, not to claim complete closure over every photonic observable.

\subsection{Training and inference protocol}

The forward surrogate is trained on this simulation-grounded dataset to predict the full response tensor from a candidate structure. In the current implementation, the model is optimized with an $L_1$ reconstruction objective in normalized response space while physical-space error is tracked during validation. Data augmentation is limited to symmetry-preserving flips that leave the optical response invariant under the adopted conventions. The conditional diffusion model is then trained on the same response tensor, together with auxiliary terms that encourage physical consistency and binarization during sampling. These choices matter less as architectural details than as indicators of intent: the generated structures are expected not only to resemble training samples, but to remain meaningful candidates for verified electromagnetic evaluation.

Inference follows a candidate-based protocol rather than a single-shot prediction protocol. For a given target response, the diffusion model produces multiple candidate structures. The forward surrogate is used to screen or approximately rank these candidates, after which selected structures are re-evaluated with RCWA. Final ranking uses task-aware response metrics derived from the verified field. This candidate-level procedure makes the one-to-many character of the inverse problem explicit. A useful inverse designer should not merely return one plausible structure; it should populate the feasible set densely enough that post-generation physical selection becomes worthwhile.

\subsection{Tasks, metrics, and scope of evaluation}

The large-NA, high-transmission, broadband second-order differentiator is used as the representative demonstration because it demands a highly structured angular response while remaining physically interpretable in both response space and image space. Additional tasks are chosen not to enlarge an application portfolio, but to test whether the same response-space formulation remains coherent under controlled deformations of the target. Fourth-order differentiation changes the target curvature. Low-pass and high-pass tasks change the target envelope. Polarization-independent and polarization-multiplexed tasks probe whether polarization enters naturally into the same formulation. A spatiotemporal-window target is retained as an extended example of a broader response-domain objective.

The evaluation metrics mirror this organization. For row-wise operator targets, the ranking logic separates center suppression, target-shape agreement, edge behavior, and effective bandwidth. For polarization-sensitive tasks, the analysis adds joint scores, passive-channel suppression, and channel matching around $\pm40^\circ$. In the baseline evaluation utilities, best-case score, mean score, top-3 score, success rate above threshold, spectrum MAE, binary rate, and candidate diversity are all computed explicitly. These quantities are not reported as an arbitrary collection of engineering diagnostics. They correspond to the central questions of the paper: whether the response shape is correct, whether physical verification agrees with nominal prediction, whether candidate diversity is preserved, and whether the generated structures remain practically usable as binary metasurfaces.
